% =====================================================================
% P1-15 -- First full manuscript draft (LaTeX source).
% Content: frozen theorem core T1-T6, C1-C6 (corrected C5), proved in
% P1_04/P1_05/P1_06 and independently audited in P1_14.  The candidate
% corollaries N1-N3 of P1-14 are NOT used here; they are collected in
% P1_15_OPTIONAL_COROLLARIES.md pending author approval.
% No novelty or priority claim is made anywhere in this draft.
% P1-17 (21 Oct 2026): audit corrections of P1_16_INDEPENDENT_AUDIT.md
% applied; see P1_17_CORRECTION_LOG.md. Core statements unchanged.
% Compilation: pdflatex P1_15_MANUSCRIPT_DRAFT.tex (twice).  Only
% standard packages (amsmath, amssymb, amsthm, geometry) are required.
% =====================================================================
\documentclass[10pt,journal,draftclsnofoot,onecolumn]{IEEEtran}

% P1-41: IEEE Transactions on Information Theory single-column review
% version, converted from P1_15_MANUSCRIPT_DRAFT.tex (commit 3e4f87f).
% Content, statements, proofs, examples, tables and references are
% unchanged; only the class/format was converted (IEEEtran review
% options above). The article-class file remains the reading copy.
\usepackage{amsmath,amssymb,amsthm}

\newtheorem{theorem}{Theorem}[section]
\newtheorem{lemma}[theorem]{Lemma}
\newtheorem{proposition}[theorem]{Proposition}
\newtheorem{corollary}[theorem]{Corollary}
\theoremstyle{definition}
\newtheorem{definition}[theorem]{Definition}
\newtheorem{example}[theorem]{Example}
\theoremstyle{remark}
\newtheorem*{remark}{Remark}

\newcommand{\Fq}{\mathbb{F}_q}
\newcommand{\Fp}{\mathbb{F}_p}
\newcommand{\Z}{\mathbb{Z}}
\newcommand{\Hull}{\operatorname{Hull}}
\newcommand{\lcm}{\operatorname{lcm}}
\newcommand{\tr}{\operatorname{tr}}
\newcommand{\E}{\mathbb{E}}

\title{The joint code--hull distribution of repeated-root constacyclic codes
over finite products of finite fields via transfer matrices}
\author{Arun Kumar Yadav\thanks{The author is with the School of Basic and Applied Sciences, K.R. Mangalam University, Gurugram, India (e-mail: arunyadav0307@gmail.com).}}

\begin{document}
\maketitle

\begin{abstract}
We study \(\lambda\)-constacyclic codes of length \(n = P n_0\), with
\(P\) a power of the characteristic \(p\) and \(\gcd(n_0, p) = 1\), over a
finite product \(A \cong \prod_s K_s\) of finite fields, for twists
compatible with the \(k\)-Galois pairing. Such codes are parametrized by
exponent vectors on the irreducible factors of \(x^{n_0} - \mu_s\). We
show that the Galois dual and hull are governed by the reciprocal
involution on those factors, and that the hull dimension is a sum of
nearest-neighbour edge terms \(\min(P - \alpha_i, \alpha_{i-1})\) along
its orbits. This gives a transfer-matrix formula for the joint
distribution of the code dimension and the hull dimension over all
labeled exponent assignments; the total number of codes; the number of
LCD codes, \(2^{\#\mathrm{orbits}}\), independent of \(P\); counts of
self-orthogonal and dual-containing codes; and the exact mean and
variance of the hull dimension, with orbits of length one treated
separately. Worked examples are verified by explicit construction of
codes.
\end{abstract}

\begin{IEEEkeywords}
Constacyclic codes, repeated-root codes, Galois hull, joint enumerator, transfer matrix, finite fields.
\end{IEEEkeywords}

\noindent Mathematics Subject Classification: 94B15, 94B05, 11T71.

% =====================================================================
\section{Introduction}
% =====================================================================

For a linear code \(C\) over a finite field, the hull --- the intersection
of \(C\) with its dual --- controls several constructions, including the
amount of entanglement needed in entanglement-assisted quantum codes and
the behaviour of LCD (linear complementary dual) codes. Replacing the
Euclidean inner product by a Galois inner product~\cite{FZ2017} leads to the
\(k\)-Galois hull \(\Hull_k(C) = C \cap D(C)\), where \(D(C)\) is the dual
with respect to a Frobenius-twisted pairing. Galois hulls of constacyclic
codes have been determined, at the level of hull dimensions and hull
generators, in a number of settings. Over finite fields,
Debnath, Prakash and Islam determined the Galois hull dimension of
constacyclic codes and the number of such codes with a given Galois hull
dimension \cite{DPI2023} (see also the correction \cite{DPI2023corr}); average Galois hull dimensions over finite
fields and over an auxiliary ring were computed by Debnath and Prakash
\cite{DP2025}, and constacyclic codes with small Galois hull dimension
were studied by Debnath, Islam, Yadav and Prakash \cite{DIYP2026}. For general linear codes, mass formulas counting codes by prescribed hull dimension were obtained in \cite{LSL2025} (with the binary case treated in \cite{LS2024}), these counts were shown to decrease as the hull dimension increases in \cite{BBO2026}, and constructions producing linear codes with prescribed Galois hull dimensions were given in \cite{LYW2025}. Over
affine algebra rings, Debnath, Islam, Mart\'inez-Moro and Prakash obtained
Galois hull formulas for constacyclic codes \cite{DIMP2026}, working
mainly in a simple-root (idempotent) framework with a repeated-root case
under additional hypotheses on the order of the twist. Relative hulls for general linear codes, including \(e\)-Galois hulls and their quantum-code consequences, were studied in \cite{ACLMRS2024}. Eldin and Sol\'e
developed generator polynomial matrices for the Galois hulls of
multi-twisted codes, of which constacyclic codes are the index-one case
\cite{ES2026}. Galois hulls of cyclic serial codes over a finite chain ring are studied in
\cite{Talbi2021}; hulls of separable double cyclic codes over \(\mathbb{Z}_{p^2}\), including \(p\)-dimensions, average \(p\)-dimensions and counts at fixed hull \(p\)-dimension, were determined in \cite{DPDC2025}, and the existence of cyclic and negacyclic codes with one-dimensional hulls was characterized in \cite{DS2024}; Galois hulls have likewise been used for quantum-code
constructions over a finite non-chain ring \cite{ZKZ2026}; repeated-root constacyclic codes and their (Euclidean)
duals, including LCD and self-dual enumerations in special lengths, have
a substantial literature \cite{CDL2015, Sharma2015, WZLL2023, ZWLL2025}.

\paragraph{Contributions and positioning.}
In this paper we (1) determine the \(k\)-Galois dual and hull of every
repeated-root constacyclic code of length \(P n_0\) over a finite product
of finite fields, for compatible twists, in exponent form
(Theorems~\ref{thm:T3} and~\ref{thm:T4}); (2) show that the joint law of
\((\dim_{\Fq} C, \dim_{\Fq} \Hull_k(C))\) over labeled assignments is a
product of traces of \((P{+}1)\)-state transfer matrices indexed by the
orbits of the reciprocal involution (Theorem~\ref{thm:T6}), with the
symmetry \(K \leftrightarrow \Delta - K\)
(Corollary~\ref{cor:jointsym}); (3) give closed forms for the numbers of
LCD, self-orthogonal and dual-containing codes
(Corollaries~\ref{cor:C3}, \ref{cor:LCDref}, \ref{cor:C4}
and~\ref{cor:restricted}); and (4) compute the exact mean and variance
of the hull dimension under the uniform labeled model, with separate
regimes for orbit lengths \(1\), \(2\) and \(\ge 3\)
(Theorems~\ref{thm:C5} and~\ref{thm:C6}). The transfer-matrix form also
has computational content: at numerical values of \((u,z)\), each orbit
trace can be evaluated by binary powering, using \(O(\log L)\)
multiplications of \((P{+}1) \times (P{+}1)\) matrices, whereas direct
enumeration inspects all \((P{+}1)^L\) assignments of the orbit.
Table~\ref{tab:positioning} positions these results against the closest
earlier work.

\begin{table}[htbp]
\centering
\small
\caption{Positioning relative to the closest earlier work. The
comparison with \cite{DPI2023} and its correction is based on the
published abstracts; no theorem-level comparison with that work is
asserted.}
\label{tab:positioning}
\begin{tabular}{@{}p{3.0cm}p{4.4cm}p{7.0cm}@{}}
\hline
Work & Setting & What is determined \\
\hline
\cite{DP2025} & constacyclic codes over finite fields and over an
auxiliary ring & average Galois hull dimensions \\
\cite{DPI2023, DPI2023corr} & constacyclic codes over finite fields &
Galois hull dimension of a given code; number of codes with given hull
dimension (as reported in the abstract) \\
\cite{DIMP2026} & constacyclic codes over affine algebra rings, mainly
in a simple-root framework, with a repeated-root case under additional
hypotheses on the order of the twist & Galois hull formulas for a given
code; entanglement-assisted constructions \\
\cite{ES2026} & multi-twisted codes, of which constacyclic codes are
the index-one case & Galois hull via generator polynomial matrices \\
This paper & repeated-root length \(n = P n_0\) with compatible twists,
over a finite product of finite fields & joint distribution of
\((\dim C, \dim \Hull_k(C))\); LCD, self-orthogonal and dual-containing
counts; exact mean and variance \\
\hline
\end{tabular}
\end{table}

The present paper studies a distributional question in the repeated-root
setting. Fix a finite product of finite fields
\(A \cong \prod_{s=1}^{N} K_s\) and a length \(n = P n_0\) in which the
characteristic divides \(n\) through a prime power \(P\) of \(p\),
while \(\gcd(n_0, p) = 1\). Every constacyclic code over \(A\) is then
described, component by component, by exponents
\(a_f \in \{0, \dots, P\}\) attached to the distinct irreducible factors
of a separable polynomial \(x^{n_0} - \mu_s\). Rather than the dimension of
a single hull, we determine the \emph{joint distribution} of the pair
\((\dim_{\Fq} C, \dim_{\Fq} \Hull_k(C))\) over all labeled exponent
assignments, for compatible twists. The joint law is captured by a
transfer-matrix enumerator: each orbit of the reciprocal involution on
the factor set contributes the trace of a power of a \((P{+}1)\)-state
matrix whose entries record, edge by edge, the code-dimension and
hull-dimension contributions (Theorem~\ref{thm:T6}). From the enumerator
we extract the total number of codes, the code-dimension distribution,
the number of LCD codes, characterizations and counts of self-orthogonal
and dual-containing codes, and exact first and second moments of the hull
dimension. Orbits of length one enter through a separate constant:
their hull edge is a self-loop \(\min(P-A, A)\), with mean
\(\nu_P = \lfloor P^2/4 \rfloor / (P+1)\), rather than the ordinary-edge
mean \(\mu_P = P(2P+1)/(6(P+1))\); the two constants must not be
interchanged when a fixed factor of the reciprocal involution is
present.

The paper is organized as follows. Section~\ref{sec:prelim} fixes the
conventions, in particular the code-first Galois pairing on which the
dual formulas depend. Sections~\ref{sec:param}--\ref{sec:hull} develop
the parametrization, the Galois dual, and the hull formula in exponent
form. Section~\ref{sec:enumerator} proves the transfer-matrix enumerator
and its specialization at \(P=1\). Sections~\ref{sec:corollaries} and
\ref{sec:moments} contain the counting corollaries and the moment
formulas. Section~\ref{sec:examples} gives worked examples and describes
the computational verification; Section~\ref{sec:limitations} records
the scope and limitations. Throughout, counts are over labeled exponent
assignments, not over equivalence or isometry classes of codes.

% =====================================================================
\section{Preliminaries and conventions}\label{sec:prelim}
% =====================================================================

Let \(p\) be a prime, \(q = p^e\), and let \(0 \le k < e\) be a fixed
Galois index. Let
\[
A \;\cong\; \prod_{s=1}^{N} K_s, \qquad
K_s = \mathbb{F}_{q^{m_s}} = \mathbb{F}_{p^{d_s}}, \quad d_s = e m_s,
\]
with fixed component labels. (Standard terminology for constacyclic codes
follows \cite{HP2003, MS1977}.)

\emph{Terminology.} We call \(A\) an \emph{affine algebra}: equivalently,
\(A\) is a finite-dimensional commutative \(\Fq\)-algebra of the form
\(\Fq[X_1,\dots,X_\ell]/\langle t_1(X_1),\dots,t_\ell(X_\ell)\rangle\)
in which each \(t_i\) is \emph{square-free}, i.e.\ has no repeated
irreducible factor; square-freeness is what makes \(A\) reduced
(\'etale), hence a finite product of finite fields. Only this
decomposition, the component Galois twists, and the weights
\(w_O = m_s \deg f\) enter the results below, so we state the setting
throughout as a finite product of finite fields. A \emph{labeled}
assignment is an exponent tuple together with the fixed component and
factor labels; all counts are over labeled assignments. On \(K_s\) write
\(\sigma(\alpha) = \alpha^{p^k}\) and
\(\rho(\alpha) = \alpha^{p^{d_s - k}} = \sigma^{-1}\).

\begin{definition}[code-first Galois pairing]
For one component (the subscript \(s\) suppressed), the \(k\)-Galois
pairing on \(K^n\) is
\(\langle x, y \rangle_k = \sum_{i=0}^{n-1} x_i \, \sigma(y_i)\),
the \emph{code-first} dual of a code \(C \subseteq K^n\) is
\(D(C) = \{ y \in K^n : \langle c, y \rangle_k = 0 \text{ for all } c \in C\}\),
and the \(k\)-Galois hull is \(\Hull_k(C) = C \cap D(C)\).
The Euclidean dual is denoted \(C^\perp\).
\end{definition}

This is the \(k\)-Galois inner product introduced by Fan and
Zhang~\cite{FZ2017}. The code-first slot order agrees with the Galois
dual convention of Liu and Pan~\cite{LP2020}; in the finite-field case
the dictionary is \(D(C) = \sigma^{e-k}(C^\perp)\). The slot order is part of the definition: the candidate \(y\) occupies
the \(\sigma\)-twisted slot, and this is what makes
\(D(C) = \rho(C^\perp)\) (coordinatewise \(\rho\);
Proposition~\ref{prop:transport}) rather than the analogous statement
with \(\sigma\). All dual formulas below use this convention. The global
pairing on \(A^n\) is the tuple of component pairings; a code over \(A\)
is an \(A\)-submodule of \(A^n\).

Fix a power \(P\) of \(p\), and \(n = P n_0\) with
\(\gcd(n_0, p) = 1\). Let \(\lambda = (\lambda_s) \in A^*\) and
\(M_s(x) = x^n - \lambda_s \in K_s[x]\). Since \(a \mapsto a^P\) is a
bijection of the finite field \(K_s\), there is a unique \(\mu_s \in K_s^*\)
with \(\mu_s^P = \lambda_s\). Writing \(\bar M_s(x) = x^{n_0} - \mu_s\), we
have \(M_s = \bar M_s^{\,P}\) in characteristic \(p\), and \(\bar M_s\) is
separable because \(\bar M_s' = n_0 x^{n_0-1}\) and every root of
\(\bar M_s\) is nonzero. Let \(\mathcal{F}_s\) be the set of distinct
monic irreducible factors of \(\bar M_s\) over \(K_s\); then
\(M_s = \prod_{f \in \mathcal{F}_s} f^P\).

An \emph{exponent assignment} on component \(s\) is a tuple
\(a = (a_f)_{f \in \mathcal{F}_s}\) with \(a_f \in \{0, \dots, P\}\); put
\(g_a = \prod_f f^{a_f}\) and \(C_s(a) = \langle g_a \rangle \subseteq
K_s[x]/(M_s)\). A global assignment is a tuple of component assignments.

A \emph{\(\lambda\)-constacyclic code} over \(A\) is an \(A\)-submodule of
\(A^n\) closed under the componentwise \(\lambda\)-constacyclic shift
\((c_0,\dots,c_{n-1}) \mapsto (\lambda c_{n-1}, c_0, \dots, c_{n-2})\).
Under the coefficient identification \(A^n \cong A[x]/(x^n-\lambda)\), this
shift is multiplication by \(x\); hence the \(\lambda\)-constacyclic codes
over \(A\) are exactly the ideals of \(A[x]/(x^n-\lambda)\), an equivalence
that needs no separability or compatibility hypothesis. From now on,
``code over \(A\)'' means a \(\lambda\)-constacyclic code over \(A\);
componentwise these are tuples of \(\lambda_s\)-constacyclic codes
(Proposition~\ref{prop:assembly}).

For a monic polynomial \(f(x) = \sum_{i=0}^{r} c_i x^i\) with
\(f(0) = c_0 \ne 0\), define the (Frobenius-twisted) reciprocal
\[
f^{\#}(x) \;=\; \rho(c_0)^{-1} \sum_{i=0}^{r} \rho(c_i)\, x^{\,r-i},
\]
which is monic of degree \(r\). Every \(f \in \mathcal{F}_s\) satisfies
\(f(0) \ne 0\) since \(\bar M_s(0) = -\mu_s \ne 0\).

\begin{definition}[compatibility and the involution \(\tau_s\)]
The twist \(\lambda_s\) is \emph{compatible} if
\(\lambda_s^{\,1 + p^{d_s - k}} = 1\). Since \(K_s^*\) is cyclic, the
number of compatible twists in \(K_s^*\) is
\(\gcd\big(1 + p^{d_s-k},\, p^{d_s} - 1\big)\); in particular
\(\lambda_s = 1\) is always compatible, and \(\lambda_s = -1\) is
compatible when \(p\) is odd.
Under compatibility, \(\tau_s : \mathcal{F}_s \to \mathcal{F}_s\),
\(\tau_s(f) = f^{\#}\), is a permutation (Theorem~\ref{thm:T2}). We
assume compatibility for every \(s\) from Section~\ref{sec:dual} on.
A \(\tau_s\)-orbit is written \(O = (f_0, \dots, f_{L-1})\) with
\(\tau_s(f_i) = f_{i+1}\), indices modulo \(L = L_O\); we set
\(\alpha_i = a_{f_i}\) and call \(\alpha_{i-1}\) the predecessor of
\(\alpha_i\). Since \(\#\) preserves degrees, all factors in an orbit
have the same degree; the \emph{weight} of \(O\) and of a factor \(f\) is
\(w_O = m_s \deg_{K_s} f\) and \(w_f = m_s \deg_{K_s} f\).
\end{definition}

Dimensions over \(\Fq\) are recovered from dimensions over \(K_s\) by
\(\dim_{\Fq} V = m_s \dim_{K_s} V\). We write \(u, z\) for indeterminates
and work in \(\Z[u,z]\).

% =====================================================================
\section{Parametrization and the reciprocal operation}\label{sec:param}
% =====================================================================

\begin{theorem}[parametrization]\label{thm:T1}
The map \(a \mapsto C_s(a) = \langle g_a \rangle\) is a bijection from
\(\{0,\dots,P\}^{\mathcal{F}_s}\) onto the set of \(\lambda_s\)-constacyclic
codes of length \(n\) over \(K_s\), and
\(\dim_{K_s} C_s(a) = n - \sum_{f} a_f \deg f\).
Consequently there are \((P+1)^{\sum_s |\mathcal{F}_s|}\) labeled global
assignments, and \(\lambda\)-constacyclic codes over \(A\) correspond to tuples
of component assignments (Proposition~\ref{prop:assembly}).
\end{theorem}

\begin{proof}
Ideals of \(K_s[x]/(M_s)\) correspond to ideals \((M_s) \subseteq I\) of
the principal ideal domain \(K_s[x]\), hence to monic divisors
\(\bar g\) of \(M_s\); by unique factorization and \(M_s =
\prod f^P\), monic divisors are exactly the \(g_a\). For a monic divisor
\(g\) of degree \(r\), multiplication by \(g\) identifies
\(\langle g \rangle \subseteq K_s[x]/(M_s)\) with
\(K_s[x]/((M_s)/g)\), of dimension \(n - r\). No separability hypothesis
is used.
\end{proof}

\begin{theorem}[the reciprocal operation]\label{thm:T2}
\begin{enumerate}
\item[(i)] \(\#\) is multiplicative on monic polynomials with nonzero
constant term; hence
\(\big(\prod_f f^{a_f}\big)^{\#} = \prod_f (f^{\#})^{a_f}\) for all
exponents \(a_f \ge 0\).
\item[(ii)] If \(f(\xi) = 0\) with \(\xi \ne 0\) in an algebraic
closure, then \(f^{\#}\big(\xi^{-p^{d_s-k}}\big) = 0\); in fact
\(f^{\#}\) is the minimal polynomial of \(\xi^{-p^{d_s-k}}\).
\item[(iii)] \(\lambda_s^{1+p^{d_s-k}} = 1 \iff
\mu_s^{1+p^{d_s-k}} = 1\). Under compatibility, \(\tau_s\) permutes
\(\mathcal{F}_s\) and preserves degrees.
\end{enumerate}
\end{theorem}

\begin{proof}
Let \(\bar\rho\) be the coefficientwise extension of \(\rho\) to
\(K_s[x]\); it is a ring automorphism with inverse \(\bar\sigma\).
For (i), \(\overline\rho(fg) = \bar\rho(f)\bar\rho(g)\), reversal
\(x^{r}f(1/x)\) is multiplicative in fixed degrees, and the normalizing
constants multiply since \((fg)_0 = f_0 g_0\); combining the three gives
\((fg)^{\#} = f^{\#} g^{\#}\), and induction on the number of factors
gives the displayed identity. For (ii), apply the Frobenius map
\(t \mapsto t^{p^{d_s-k}}\) of the algebraic closure to
\(\sum c_i \xi^i = 0\) to get \(\sum \rho(c_i) \beta^{-i} = 0\) for
\(\beta = \xi^{-p^{d_s-k}}\); multiplying by \(\beta^r\) yields
\(\rho(c_0) f^{\#}(\beta) = 0\). The same Frobenius transport preserves
degrees of minimal polynomials, so \(f^{\#}\) is the (monic irreducible)
minimal polynomial of \(\beta\). For (iii),
\((\mu_s^{1+p^{d_s-k}})^P = \lambda_s^{1+p^{d_s-k}}\) and in
characteristic \(p\), \(y^P = 1\) implies \(y = 1\); now (ii) shows
\(\tau_s(f)\) is the minimal polynomial of a root \(\beta\) with
\(\beta^{n_0} = \rho(\mu_s)^{-1} = \mu_s\) (using compatibility in the
form \(\rho(\mu_s) = \mu_s^{-1}\)), so \(\tau_s(f) \mid \bar M_s\);
bijectivity follows by applying the same construction with \(\sigma\) in
place of \(\rho\), which gives the inverse map.
\end{proof}

% =====================================================================
\section{The Galois dual}\label{sec:dual}
% =====================================================================

The first step is a Euclidean-dual statement that uses no separability
at all.

\begin{lemma}[Euclidean dual generator]\label{lem:RR3}
Let \(g\) be a monic divisor of \(x^n - \lambda\), \(\lambda \in K^*\),
\(h = (x^n-\lambda)/g\), \(r = \deg g\), \(s = n - r\), and
\(C = \langle g \rangle\) under the coefficient identification (this is
the classical constacyclic dual statement;
cf.~\cite{HP2003, MS1977}). Then
\(h_0 \ne 0\), and with
\(h^{\natural}(x) = h_0^{-1} \sum_{i=0}^{s} h_i x^{s-i}\) (the ordinary
reciprocal, without Frobenius),
\(C^\perp = \langle h^{\natural} \rangle\), a \(\lambda^{-1}\)-constacyclic
code.
\end{lemma}

\begin{proof}
From \(gh = x^n - \lambda\), \(g_0 h_0 = -\lambda \ne 0\). Both
\(\langle h^{\natural} \rangle\) and \(C^\perp\) have dimension \(r\):
the first by the divisor dimension count of Theorem~\ref{thm:T1}
(applied to any modulus), the second by rank--nullity. The boundary
cases are immediate from the definitions: if \(r = 0\), then \(g = 1\),
\(C = K^n\) and \(C^\perp = \{0\}\), while \(h^{\natural}\) has degree
\(n\) and hence no nonzero multiple of degree \(< n\); if \(s = 0\), then
\(h = 1\), \(C = \{0\}\) and \(C^\perp = K^n\), while
\(h^{\natural} = 1\). Assume \(r, s \ge 1\) below, so that \(V\) is
defined (and \(v = 0\) gives \(d = 0\)). We show orthogonality. For polynomials \(a, d\) of degree \(< n\), write
\(\operatorname{rev}_n(d) = x^{n-1} d(1/x)\); then
\(\sum_i a_i d_i = [x^{n-1}]\big(a(x) \operatorname{rev}_n(d)(x)\big)\).
Take \(a = gu\), \(\deg u < s\), and \(d = h^{\natural} v\),
\(\deg v < r\). Since
\(h^{\natural}(1/x) = h_0^{-1} x^{-s} h(x)\), one computes
\(\operatorname{rev}_n(d) = h_0^{-1} h(x) V(x)\) with
\(V(x) = x^{r-1} v(1/x)\), \(\deg V \le r-1\). Hence
\(a(x)\operatorname{rev}_n(d)(x) = h_0^{-1} (x^n - \lambda) u(x) V(x)\)
with \(\deg(uV) \le n-2\); the coefficient of \(x^{n-1}\) is \(0\).
Thus \(\langle h^{\natural}\rangle \subseteq C^\perp\), and equality
follows by dimensions. That \(h^{\natural} \mid x^n - \lambda^{-1}\)
follows from \((gh)^* = g^* h^* = -\lambda(x^n - \lambda^{-1})\) for the
unnormalized reciprocal \(\phi^*(x) = x^{\deg\phi}\phi(1/x)\).
\end{proof}

\begin{proposition}[transport by \(\rho\)]\label{prop:transport}
Let the field be \(\mathbb{F}_{p^d}\), with \(\sigma, \rho\) as in
Section~\ref{sec:prelim}. For every \(\mathbb{F}_{p^d}\)-linear code
\(C \subseteq \mathbb{F}_{p^d}^n\),
\(D(C) = \rho(C^\perp)\) (coordinatewise). Consequently, for
\(C = \langle g \rangle \subseteq \mathbb{F}_{p^d}[x]/(x^n-\lambda)\) and
\(h = (x^n-\lambda)/g\),
\(D(C) = \langle h^{\#} \rangle \subseteq
\mathbb{F}_{p^d}[x]/(x^n - \lambda^{-p^{d-k}})\); if \(\lambda\) is
compatible, the modulus is again \(x^n - \lambda\).
\end{proposition}

\begin{proof}
By the code-first definition, \(y \in D(C)\) iff
\(\sum_i c_i \sigma(y_i) = 0\) for all \(c \in C\), i.e.\ iff
\(\sigma(y) \in C^\perp\) (coordinatewise), i.e.\ iff
\(y \in \rho(C^\perp)\) since \(\rho = \sigma^{-1}\); this is exactly
where the slot order enters. Coefficientwise application of an
automorphism sends the ideal \(\langle e \rangle\) in
\(K[x]/(x^n-\eta)\) to \(\langle \bar\rho(e)\rangle\) in
\(K[x]/(x^n - \rho(\eta))\); applying this to Lemma~\ref{lem:RR3} and
comparing \(\bar\rho(h^{\natural})\) with \(h^{\#}\) gives the generator.
Compatibility \(\lambda^{1+p^{d-k}}=1\) gives
\(\lambda^{-p^{d-k}} = \lambda\).
\end{proof}

\begin{theorem}[dual exponents]\label{thm:T3}
For every assignment \(a\), with \(h_a = M_s / g_a\),
\(D(C_s(a)) = \langle h_a^{\#} \rangle\); under compatibility this is a
\(\lambda_s\)-constacyclic code in the same quotient. On a
\(\tau_s\)-orbit \(O = (f_0,\dots,f_{L-1})\) with \(\alpha_i = a_{f_i}\),
the exponent of \(f_i\) in \(h_a^{\#}\) is
\[
\beta_i \;=\; P - \alpha_{i-1} \qquad (\text{indices mod } L).
\]
\end{theorem}

\begin{proof}
The generator statement is Proposition~\ref{prop:transport}. For the
exponents, \(h_a = \prod_{f} f^{P - a_f}\), so by Theorem~\ref{thm:T2}(i)
and the permutation property,
\(h_a^{\#} = \prod_f (\tau_s f)^{P-a_f} = \prod_{g} g^{P - a_{\tau_s^{-1} g}}\);
for \(g = f_i\), \(\tau_s^{-1}(f_i) = f_{i-1}\), giving \(P - \alpha_{i-1}\).
\end{proof}

\begin{remark}
Without compatibility, \(D(C_s(a)) = \langle h_a^{\#}\rangle\) and the
exponent formula remain valid, but in the modulus
\(x^n - \lambda_s^{-p^{d_s-k}}\). Compatibility is what allows the dual
to be viewed inside the same constacyclic family; it cannot simply be
dropped (e.g.\ in \(\mathbb{F}_9\) with \(k=1\), a twist \(\lambda\) of
order \(8\) satisfies \(\lambda^{1+3} = \lambda^4 = -1 \ne 1\) and indeed
\(\lambda^{-3} \ne \lambda\)).
\end{remark}

% =====================================================================
\section{The Galois hull}\label{sec:hull}
% =====================================================================

\begin{lemma}[intersection of ideals]\label{lem:lcm}
For monic divisors \(g, l\) of a monic \(M \in K[x]\), the images of
\((g)\) and \((l)\) in \(K[x]/(M)\) intersect in the image of
\((\lcm(g,l))\), and
\(\lcm(g,l) = \prod_{\pi} \pi^{\max(v_\pi(g), v_\pi(l))}\).
\end{lemma}

\begin{proof}
In \(K[x]\), \((g) \cap (l) = (\lcm(g,l))\) since both consist of the
common multiples of \(g\) and \(l\), characterized valuation-wise by the
displayed maximum. If \(y\) lies in both images in the quotient, any
lift \(Y\) satisfies \(Y \in (g)\) and \(Y \in (l)\) because
\((M) \subseteq (g), (l)\); hence \(Y \in (\lcm(g,l))\). The reverse
inclusion is immediate from \(g, l \mid \lcm(g,l)\).
\end{proof}

\begin{theorem}[hull generator and dimensions]\label{thm:T4}
Under compatibility, for every assignment \(a\),
\(\Hull_k(C_s(a)) = \langle \lcm(g_a, h_a^{\#}) \rangle\), the exponent
of \(f_i\) in the hull generator being \(\max(\alpha_i, P-\alpha_{i-1})\).
Moreover
\[
\dim_{\Fq} C(a) \;=\; \sum_{s} \sum_{O} w_O \sum_{i} (P - \alpha_i),
\qquad
\dim_{\Fq} \Hull_k(C(a)) \;=\; \sum_{s} \sum_{O} w_O \sum_i
\min(P-\alpha_i,\, \alpha_{i-1}),
\]
the inner sums over orbit positions (indices mod \(L_O\)); here and
below, \(\alpha_i\) refers to the orbit \((s,O)\) under discussion.
\end{theorem}

\begin{proof}
Work in one component. By Theorems~\ref{thm:T1} and \ref{thm:T3},
\(C(a)\) and \(D(C(a))\) are the images of \((g_a)\) and
\((h_a^{\#})\), both dividing \(M = x^n-\lambda\); Lemma~\ref{lem:lcm}
gives the hull generator, whose \(f_i\)-exponent is
\(\max(\alpha_i, \beta_i) = \max(\alpha_i, P-\alpha_{i-1})\). Dimensions
are \(n - \deg(\text{generator})\) over \(K\); grouping factors by
orbits (constant degree on an orbit) and using
\(P - \max(u,v) = \min(P-u, P-v)\) with \(v = P-\alpha_{i-1}\) gives the
hull formula; multiplication by \(m_s\) and summation over \(s\) gives
the \(\Fq\)-dimensions.
\end{proof}

% =====================================================================
\section{The transfer-matrix enumerator}\label{sec:enumerator}
% =====================================================================

For \(P \ge 1\), \(w \ge 1\), define the \((P{+}1)\times(P{+}1)\) matrix
over \(\Z[u,z]\)
\[
T^{(P)}_w(u,z)_{a,b} \;=\; u^{\,w(P-b)}\, z^{\,w \min(P-b,\, a)},
\qquad a, b \in \{0,\dots,P\},
\]
the row index \(a\) being the previous exponent and the column index
\(b\) the current one.

\begin{theorem}[local transfer trace]\label{thm:T5}
For a \(\tau_s\)-orbit \(O\) of length \(L\) and weight \(w\),
\[
\tr\!\big(T^{(P)}_w(u,z)^L\big)
\;=\; \sum_{\alpha \in \{0,\dots,P\}^L}
u^{\,w \sum_i (P-\alpha_i)}\,
z^{\,w \sum_i \min(P-\alpha_i,\, \alpha_{i-1})},
\quad \text{indices mod } L.
\]
\end{theorem}

\begin{proof}
Expanding the matrix power,
\((T^L)_{a_0,a_L} = \sum_{a_1,\dots,a_{L-1}} \prod_{j=1}^{L} T_{a_{j-1},a_j}\);
setting \(a_L = a_0\) and summing over \(a_0\) gives
\(\tr(T^L) = \sum_{a_0,\dots,a_{L-1}} \prod_{j=1}^{L} T_{a_{j-1},a_j}\).
These closed walks are in bijection with assignments
\(\alpha = (\alpha_0,\dots,\alpha_{L-1})\) (orbit positions are labeled;
a walk and its rotation are different assignments). The \(j\)-th edge
contributes \(u^{w(P-a_j)} z^{w\min(P-a_j, a_{j-1})}\); summing exponents
over the walk and reindexing gives the claim. For \(L=1\) the single edge
is the self-loop \(a_0 \to a_0\) and the formula reads
\(\sum_a u^{w(P-a)} z^{w\min(P-a,a)}\), as required.
\end{proof}

\begin{proposition}[product assembly]\label{prop:assembly}
Ideals of \(A[x]/(x^n-\lambda) \cong \prod_s K_s[x]/(x^n-\lambda_s)\) are
exactly the products of component ideals; the global dual and hull
decompose componentwise; and
\(\dim_{\Fq} C = \sum_s m_s \dim_{K_s} C_s\), likewise for the hull.
\end{proposition}

\begin{proof}
The ring isomorphism is coefficientwise; quotienting a product by the
product ideal \((x^n-\lambda_s)_s\) gives the product of quotients (this
is the direct-product structure of the square-free algebra \(A\) itself;
no coprimality between components is involved). Ideals of a finite
product of rings are products of ideals (project, and recover the
components with the idempotents of \(A\)). The pairing is componentwise
and vanishes exactly when every component does, giving the dual and hull
statements; dimensions add over the direct sum \(A^n \cong \bigoplus_s
K_s^n\) with the factor \([K_s:\Fq] = m_s\).
\end{proof}

\begin{theorem}[joint enumerator]\label{thm:T6}
Under the standing hypotheses (common \(P, n, k\); compatible twists;
labeled components),
\[
\mathcal{E}(u,z) \;:=\;
\sum_{\text{labeled } a} u^{\dim_{\Fq} C(a)} z^{\dim_{\Fq} \Hull_k(C(a))}
\;=\; \prod_{s=1}^{N} \prod_{O \in \mathcal{O}_s}
\tr\!\Big(T^{(P)}_{w_O}(u,z)^{L_O}\Big),
\]
where \(\mathcal{O}_s\) is the set of \(\tau_s\)-orbits.
The coefficient of \(u^K z^H\) is the number of labeled assignments with
\(\dim_{\Fq} C = K\) and \(\dim_{\Fq} \Hull_k(C) = H\).
\end{theorem}

\begin{proof}
By Theorem~\ref{thm:T1} and Proposition~\ref{prop:assembly}, a global
assignment is the disjoint union of its orbit restrictions
\(\alpha^{(O)} \in \{0,\dots,P\}^{L_O}\), and by Theorem~\ref{thm:T4}
each of the two dimensions is a sum over \((s,O)\) of a quantity
depending only on \(\alpha^{(O)}\). The finite distributive law
(expanding a product of sums over a Cartesian product index set)
factorizes the global sum into the product of the orbit sums
\(\mathcal{E}_O(u,z) = \sum_{\alpha} u^{w_O \sum (P-\alpha_i)}
z^{w_O \sum \min(P-\alpha_i,\alpha_{i-1})}\), and Theorem~\ref{thm:T5}
identifies \(\mathcal{E}_O(u,z) = \tr(T^{(P)}_{w_O}(u,z)^{L_O})\).
\end{proof}

\begin{corollary}[joint symmetry of the enumerator]\label{cor:jointsym}
Put \(\Delta := \sum_{s,O} P w_O L_O = n \sum_s m_s\) (indeed, for each
component \(s\), \(\sum_{O \in \mathcal{O}_s} w_O L_O = m_s n_0\), since
the degrees of the distinct irreducible factors of
\(\bar M_s\) sum to \(n_0\); hence \(\Delta\) is the maximum, and ambient,
\(\Fq\)-dimension). Then for all integers \(K, H\),
\begin{align*}
\#\{a : (\dim_{\Fq} C(a), \dim_{\Fq} \Hull_k(C(a))) = (K,H)\}
&= \#\{a : (\dim_{\Fq} C(a), \dim_{\Fq} \Hull_k(C(a))) = (\Delta-K,H)\},
\end{align*}
equivalently \(\mathcal{E}(u,z) = u^{\Delta} \mathcal{E}(u^{-1}, z)\).
In particular, for every \(K\) the multiset of hull dimensions over the
assignments with \(\dim_{\Fq} C = K\) equals that over the assignments
with \(\dim_{\Fq} C = \Delta - K\).
\end{corollary}

\begin{proof}
It suffices to argue on one orbit \(O\) of length \(L\) and weight \(w\);
the global statement follows by summing over orbits. Define
\(\alpha'_i = P - \alpha_{-i}\) (indices mod \(L\)). This is an
involution of \(\{0,\dots,P\}^L\): applying it twice returns
\(\alpha_i\). Its code contribution is
\(w \sum_i (P-\alpha'_i) = w \sum_i \alpha_{-i} = w \sum_i \alpha_i
= PwL - w\sum_i (P-\alpha_i)\), so \(K\) maps to \(PwL - K\). Its hull
contribution satisfies, with the reindexing \(j = -i+1\),
\(\min(P-\alpha'_i, \alpha'_{i-1}) = \min(\alpha_{-i}, P-\alpha_{-i+1})
= \min(P-\alpha_j, \alpha_{j-1})\),
term by term equal to the hull sum of \(\alpha\). Hence the involution
pairs tuples of joint type \((K,H)\) with tuples of type
\((PwL-K, H)\), proving the per-orbit coefficient identity; the product
over all orbits (Theorem~\ref{thm:T6}) gives the claim with
\(\Delta = \sum_O Pw_OL_O\).
\end{proof}

\begin{remark}[hypotheses of the identity]
As an identity in \(\Z[u,z]\), Theorem~\ref{thm:T5} holds for every
integer \(P \ge 1\); that \(P\) is a power of \(p\) is needed
for the field construction behind Theorems~\ref{thm:T1}--\ref{thm:T4},
not for the matrix identity. The enumerator counts labeled assignments;
the trace formula is insensitive to reversing the orbit orientation.
Indeed the hull dimension of every individual assignment is unchanged
under reversal, because
\(\min(P-a,b) - \min(P-b,a) = b - a\) telescopes to \(0\) around the
orbit; hence no statistic of Sections~\ref{sec:corollaries}
and~\ref{sec:moments} can detect the orientation. The predecessor
convention \(\alpha_{i-1}\) in Theorem~\ref{thm:T4} is nevertheless
essential for identifying the hull \emph{generator}, and is fixed by
the code-first pairing of Section~\ref{sec:prelim}.
\end{remark}

\begin{proposition}[simple-root recovery at \(P=1\)]\label{prop:P1}
At \(P = 1\), assignments are binary, \(g_a\) is a
squarefree divisor, the hull edge term \(\min(1-\alpha_i, \alpha_{i-1})\)
is the indicator of the transition \(1 \to 0\) along the orbit, and
\[
T^{(1)}_w(u,z) \;=\;
\begin{pmatrix} u^w & 1 \\ u^w z^w & 1 \end{pmatrix}
\qquad (\text{rows/columns } a,b \in \{0,1\}),
\]
so Theorem~\ref{thm:T6} specializes to the binary-state simple-root
joint enumerator with the same conventions.
\end{proposition}

\begin{proof}
For \(P=1\), \(\min(1-b, a)\) is \(1\) exactly for \((a,b) = (1,0)\);
reading off the four entries gives the matrix, and \(\mu_s = \lambda_s\),
\(\bar M_s = M_s\) gives the squarefree parametrization.
\end{proof}

% =====================================================================
\section{Counting corollaries}\label{sec:corollaries}
% =====================================================================

\begin{corollary}[total count]\label{cor:C1}
\(\mathcal{E}(1,1) = (P+1)^{\sum_s |\mathcal{F}_s|}\).
\end{corollary}

\begin{proof}
At \(u=z=1\) each \(T^{(P)}_w\) is the all-ones matrix \(J\) with
\(J^L = (P+1)^{L-1} J\), so \(\tr(J^{L_O}) = (P+1)^{L_O}\); product over
orbits gives \((P+1)^{\sum_s \sum_O L_O} = (P+1)^{\sum_s |\mathcal{F}_s|}\).
\end{proof}

\begin{corollary}[code-dimension distribution]\label{cor:C2}
\(\displaystyle
\mathcal{E}(u,1) = \prod_{s} \prod_{f \in \mathcal{F}_s}
\big(1 + u^{w_f} + \dots + u^{P w_f}\big)\).
\end{corollary}

\begin{proof}
At \(z=1\) the entry \(T_{a,b} = u^{w(P-b)}\) is independent of the row,
so \(T = c\,r\) is rank one with \(r \cdot c = \sum_{j=0}^{P} u^{wj}\);
then \(\tr(T^L) = (r\cdot c)^{L-1} \tr(T) = \big(\sum_{j=0}^P u^{wj}\big)^L\).
Within an orbit all \(w_f\) equal \(w_O\), giving
\(\prod_{f \in O} (1 + u^{w_f} + \dots + u^{Pw_f})\) per orbit.
\end{proof}

\begin{corollary}[LCD count]\label{cor:C3}
The number of labeled assignments with \(\Hull_k(C(a)) = 0\) is
\(2^{\#\text{orbits}}\), independent of \(P\).
\end{corollary}

\begin{proof}
The hull dimension is a sum of nonnegative terms
\(w_O \min(P-\alpha_i, \alpha_{i-1})\); it vanishes iff every orbit
satisfies \(\min(P-\alpha_i, \alpha_{i-1}) = 0\) for all \(i\). On one
orbit, if some \(\alpha_m \notin \{0,P\}\) (so \(P \ge 2\)), then
\(\min(P-\alpha_{m+1}, \alpha_m) = 0\) with \(\alpha_m > 0\) forces
\(\alpha_{m+1} = P\), and inductively every \(\alpha_i = P\) around the
cycle, a contradiction (for \(L=1\), \(\min(P-\alpha_0,\alpha_0) = 0\)
already forces \(\alpha_0 \in \{0,P\}\)). Hence each \(\alpha_i \in
\{0,P\}\); if both values occurred, at a boundary position with
\(\alpha_i = 0\), \(\alpha_{i-1} = P\) the term would be
\(\min(P,P) = P \ge 1\). So exactly two tuples per orbit
(\(\alpha \equiv 0\) or \(\alpha \equiv P\)), independently across
orbits.
\end{proof}

\begin{corollary}[LCD refinement: the code dimension of LCD assignments]
\label{cor:LCDref}
With \(\Delta\) as in Corollary~\ref{cor:jointsym},
\[
\sum_{a \,:\, \Hull_k(C(a)) = 0} u^{\dim_{\Fq} C(a)}
\;=\; \prod_{s} \prod_{O \in \mathcal{O}_s} \big(1 + u^{P w_O L_O}\big).
\]
Consequently, under the uniform distribution on labeled assignments,
conditionally on \(\Hull_k(C(a)) = 0\) the code dimension is distributed
as \(\sum_O (P w_O L_O)\, \xi_O\) with independent
\(\xi_O \sim \mathrm{Bernoulli}(1/2)\); in particular
\(\E[\dim_{\Fq} C \mid \Hull_k(C) = 0] = \Delta/2\) and
\(\operatorname{Var}(\dim_{\Fq} C \mid \Hull_k(C) = 0)
= \frac{P^2}{4} \sum_O (w_O L_O)^2\).
\end{corollary}

\begin{proof}
By Corollary~\ref{cor:C3}, \(\Hull_k(C(a)) = 0\) if and only if every
orbit restriction is constant, \(\alpha \equiv 0\) or \(\alpha \equiv
P\); these two tuples have code contributions \(Pw_OL_O\) and \(0\)
respectively (Theorem~\ref{thm:T4}), giving the displayed product by
independence across orbits. Conditioning the uniform distribution on
this product event keeps the orbits independent with the two options
equally likely, which is the stated Bernoulli sum; its mean is
\(\tfrac12 \sum_O Pw_OL_O = \Delta/2\) and its variance is
\(\sum_O (Pw_OL_O)^2 \cdot \tfrac14\).
\end{proof}

\begin{corollary}[self-orthogonal and dual-containing codes]
\label{cor:C4}
\(C(a) \subseteq D(C(a))\) if and only if
\(\alpha_{i-1} + \alpha_i \ge P\) on every orbit edge, and then the count
is \(\prod_O \tr(B_P^{L_O})\) with \((B_P)_{a,b} = 1\) iff \(a+b \ge P\).
Dually, \(D(C(a)) \subseteq C(a)\) if and only if
\(\alpha_{i-1} + \alpha_i \le P\) on every edge, with count
\(\prod_O \tr((B'_P)^{L_O})\), \((B'_P)_{a,b} = 1\) iff \(a+b \le P\).
The two counts are equal.
\end{corollary}

\begin{proof}
Since \(\Hull = C \cap D(C) \subseteq C\), \(C \subseteq D(C)\) iff
\(\dim \Hull = \dim C\). By Theorem~\ref{thm:T4},
\(\dim C - \dim \Hull = \sum_{s,O} w_O \sum_i \max(0, P-\alpha_i-\alpha_{i-1})\)
because \((P-\alpha_i) - \min(P-\alpha_i,\alpha_{i-1}) =
\max(0, P-\alpha_i-\alpha_{i-1})\); every summand is nonnegative, so the
total vanishes iff every edge satisfies \(\alpha_{i-1}+\alpha_i \ge P\).
The admissible tuples on one orbit are the closed walks of the \(0\)-\(1\)
matrix \(B_P\), counted by \(\tr(B_P^{L_O})\) as in
Theorem~\ref{thm:T5}. For the dual side, \(\dim D(C) =
\sum w_O \sum_i \alpha_{i-1} = \sum w_O \sum_i \alpha_i\) by
Theorem~\ref{thm:T3}, and \(\alpha_{i-1} - \min(P-\alpha_i,\alpha_{i-1})
= \max(0, \alpha_i+\alpha_{i-1}-P)\); the same argument gives the
condition \(\alpha_{i-1}+\alpha_i \le P\) and the count
\(\tr((B'_P)^{L_O})\). The involution \(\alpha_i \mapsto P-\alpha_i\)
maps one edge condition to the other, so the counts agree.
\end{proof}

\begin{corollary}[restricted enumerators for the self-orthogonal and
dual-containing families]\label{cor:restricted}
For an orbit of length \(L\) and weight \(w\), define matrices over
\(\Z[t]\) and \(\Z[u,z]\) by
\[
\big(A^{\mathrm{SO}}_w(t)\big)_{a,b} =
\begin{cases} t^{\,w(P-b)}, & a + b \ge P, \\ 0, & a + b < P,
\end{cases}
\qquad
\big(A^{\mathrm{DC}}_w(u,z)\big)_{a,b} =
\begin{cases} u^{\,w(P-b)} z^{\,wa}, & a + b \le P, \\ 0, & a + b > P,
\end{cases}
\]
the row index \(a\) being the previous exponent and the column index
\(b\) the current one, as in Theorem~\ref{thm:T5}. Then
\[
\sum_{a \,:\, C(a) \subseteq D(C(a))} t^{\dim_{\Fq} C(a)}
= \prod_{s} \prod_{O \in \mathcal{O}_s}
\tr\!\Big(\big(A^{\mathrm{SO}}_{w_O}(t)\big)^{L_O}\Big),
\]
whose specialization \(t = uz\) is the joint enumerator of the
self-orthogonal family (there \(\dim_{\Fq} \Hull_k = \dim_{\Fq} C\)), and
\[
\sum_{a \,:\, D(C(a)) \subseteq C(a)}
u^{\dim_{\Fq} C(a)} z^{\dim_{\Fq} \Hull_k(C(a))}
= \prod_{s} \prod_{O \in \mathcal{O}_s}
\tr\!\Big(\big(A^{\mathrm{DC}}_{w_O}(u,z)\big)^{L_O}\Big).
\]
Writing \(N_{\mathrm{SO}}(K)\) for the number of self-orthogonal
assignments with \(\dim_{\Fq} C = K\) and \(N_{\mathrm{DC}}(K,H)\) for
the number of dual-containing assignments with
\((\dim_{\Fq} C, \dim_{\Fq} \Hull_k) = (K,H)\), one has
\(N_{\mathrm{SO}}(K) = N_{\mathrm{DC}}(\Delta - K, K)\) for all \(K\).
At \(u = z = 1\) the two enumerators reduce to the two counts of
Corollary~\ref{cor:C4}.
\end{corollary}

\begin{proof}
Both admissibility conditions are those of Corollary~\ref{cor:C4}, and
both dimensions are sums over orbits, so by the closed-walk expansion of
Theorem~\ref{thm:T5} it suffices to check one orbit. On a
self-orthogonal tuple, \(\Hull_k(C) = C\), so the joint weight is
\((uz)^{\dim_{\Fq} C}\); along an edge \(a \to b\) the code contribution
is \(w(P-b)\) and admissibility is \(a+b \ge P\), which is exactly the
weight built into \(A^{\mathrm{SO}}_w\). On a dual-containing tuple,
\(\Hull_k(C) = D(C)\), whose per-orbit dimension is
\(w \sum_i \alpha_i = w \sum_i \alpha_{i-1}\) (proof of
Corollary~\ref{cor:C4}); the edge \(a \to b\) therefore carries code
weight \(w(P-b)\) (from \(b\)) and hull weight \(wa\) (from the previous
exponent), with admissibility \(a + b \le P\), as in
\(A^{\mathrm{DC}}_w\). The trace includes the closing edge, i.e. the
orbit-boundary condition. For the symmetry, the per-orbit complement
\(\alpha_i \mapsto P - \alpha_i\) maps the self-orthogonal condition to
the dual-containing condition bijectively, sends \(\dim_{\Fq} C = K\) to
\(\dim_{\Fq} C = \Delta - K\), and sends the hull dimension to
\(\sum_O w_O \sum_i (P - \alpha_i) = K\).
\end{proof}

% =====================================================================
\section{Moments of the hull dimension}\label{sec:moments}
% =====================================================================

\emph{The model.} By Theorem~\ref{thm:T1}, labeled assignments are in
bijection with the labeled constacyclic codes of the given length,
alphabet and twist; we therefore put the uniform measure on assignments,
equivalently on the codes of this family itself (not on random linear
codes). The mean and variance below are the exact ensemble averages of
\(\dim_{\Fq}\Hull_k\) over this family. Since
\(\mu_P / P \to 1/3\) as \(P \to \infty\) (and \(\mu_1 = 1/4\)), the mean
hull dimension is a positive proportion of the ambient dimension
\(\Delta\) --- roughly \(\Delta/3\) for large \(P\), \(\Delta/4\) at
\(P = 1\), apart from the length-one-orbit correction --- so the typical
code of the family has a hull growing linearly with \(n\), in contrast
to random linear codes, whose hull dimension stays bounded in
probability~\cite{Sendrier1997}. The variance formulas quantify the
fluctuations around this behaviour and, through the entanglement cost
\(c = n - K - H\) of the conditional quantum remark in
Section~\ref{sec:limitations}, describe its typical size in this
ensemble. Comparisons with average-dimension formulas in the literature
should be made only under matching hypotheses; the results below are
stated strictly for the labeled repeated-root ensemble defined here.

Choose a labeled assignment uniformly at random; equivalently, all
exponents \(A_i\) (over all components, orbits, positions) are
independent and uniform on \(\{0,\dots,P\}\). For one orbit write
\(X_i = \min(P-A_i, A_{i-1})\) (indices mod \(L\)) and
\(S_L = \sum_i X_i\); the orbit contributes \(w_O S_{L_O}\) to
\(\dim_{\Fq}\Hull_k\) (Theorem~\ref{thm:T4}). Set
\[
\mu_P = \frac{P(2P+1)}{6(P+1)},
\qquad
\nu_P = \frac{\lfloor P^2/4 \rfloor}{P+1}.
\]

\begin{lemma}[edge moments]\label{lem:edge}
For independent \(U, V\) uniform on \(\{0,\dots,P\}\),
\(\E[\min(U,V)] = \mu_P\) and
\(\E[\min(U,V)^2] = \frac{P(P^2+P+1)}{6(P+1)}\).
For \(0 \le c \le P\), \(F(c) := \E[\min(c,U)] =
\frac{c(2P-c+1)}{2(P+1)}\).
\end{lemma}

\begin{proof}
For a nonnegative integer variable,
\(\E[X] = \sum_{t\ge1} \Pr(X \ge t)\) and
\(\E[X^2] = \sum_{t\ge1} (2t-1)\Pr(X\ge t)\).
Here \(\Pr(\min(U,V) \ge t) = ((P-t+1)/(P+1))^2\) for \(1 \le t \le P\);
summing via \(\sum_{r=1}^P r^2 = P(P+1)(2P+1)/6\) and
\(\sum r^3 = (P(P+1)/2)^2\) gives both formulas. The last claim is
\(F(c) = \sum_{t=1}^c (P-t+1)/(P+1)\).
\end{proof}

\begin{theorem}[means]\label{thm:C5}
\(\displaystyle \E[\dim_{\Fq} C] = \frac{n}{2} \sum_s m_s\), and
\[
\E[\dim_{\Fq} \Hull_k(C)] \;=\;
\nu_P \sum_{O:\, L_O = 1} w_O \;+\;
\mu_P \sum_{O:\, L_O \ge 2} w_O L_O .
\]
\end{theorem}

\begin{proof}
\(\E[P - A_f] = P/2\) and \(\sum_{f \in \mathcal{F}_s} \deg f = n_0\)
give the first formula via Theorem~\ref{thm:T4}. For the hull, linearity
over orbits reduces to \(\E[S_L]\). If \(L \ge 2\), each \(X_i\) is the
minimum of two independent uniforms (for \(L=2\), \(X_0, X_1\) are the
two ordered pairings of the same two distinct variables), so
\(\E[S_L] = L\mu_P\) by Lemma~\ref{lem:edge}. If \(L = 1\),
\(\E[S_1] = \frac{1}{P+1}\sum_{a=0}^P \min(P-a,a)\); the summands are
symmetric under \(a \mapsto P-a\), giving \(P^2/4\) for \(P = 2m\) and
\((P^2-1)/4\) for \(P = 2m+1\), i.e.\ \(\lfloor P^2/4 \rfloor\); hence
\(\E[S_1] = \nu_P\). Note \(\nu_1 = 0\): at \(P=1\) an \(L_O=1\) orbit
contributes nothing, and \(\E[\Hull] = \frac14 \sum_{L_O\ge2} w_O L_O\).
\end{proof}

For the variance, orbit contributions are independent (disjoint exponent
variables), so with \(V_L(P) := \operatorname{Var}(S_L)\),
\(\operatorname{Var}(\dim_{\Fq}\Hull_k) = \sum_{s,O} w_O^2 V_{L_O}(P)\);
note the squared weights, in contrast to the mean.

\begin{lemma}[variance pieces]\label{lem:varpieces}
\begin{enumerate}
\item[(i)] \(L=1\):
\(V_1(P) = \dfrac{P(P+2)(P^2+2P+4)}{48(P+1)^2}\) if \(P\) is even, and
\(V_1(P) = \dfrac{(P-1)(P+3)}{48}\) if \(P\) is odd.
\item[(ii)] For \(L \ge 3\), two adjacent edges satisfy
\(\operatorname{Cov}(X_i, X_{i+1}) = \kappa_P < 0\), where
\(\kappa_P = -\frac{P(P+2)(7P^2+14P+9)}{360(P+1)^2}\), while non-adjacent edge pairs
are independent.
\item[(iii)] For \(L=2\):
\(\E[X_0 X_1] = \dfrac{P^2(P-1)}{12(P+1)}\) and
\(\operatorname{Cov}(X_0,X_1) = -\dfrac{P^2(P+2)^2}{36(P+1)^2}\).
\end{enumerate}
\end{lemma}

\begin{proof}
(i) \(\E[\min(P-A,A)^2] = \frac{1}{P+1}\sum_a \min(P-a,a)^2\); for
\(P = 2m\) the squared values are \(0^2,\dots,(m-1)^2,m^2,(m-1)^2,\dots,0^2\),
summing to \(m(2m^2+1)/3\), and for \(P=2m+1\) each of \(0^2,\dots,m^2\)
occurs twice, summing to \(m(m+1)(2m+1)/3\); subtracting \(\nu_P^2\)
(Theorem~\ref{thm:C5}) and simplifying gives the claim.
(ii) \(X_i, X_{i+1}\) share exactly the variable \(A_i\). Conditioning on
\(A_i = b\) makes them conditionally independent with means
\(F(P-b)\) and \(F(b)\) (Lemma~\ref{lem:edge}), so
\(\E[X_iX_{i+1}] = \frac{1}{P+1}\sum_{b=0}^P F(P-b)F(b)
= \frac{P(P-1)(11P^2+15P+6)}{120(P+1)^2}\) by the power sums
\(\sum b^j\), \(j \le 4\); subtracting \(\mu_P^2\) gives \(\kappa_P\),
negative for \(P \ge 1\). Non-adjacent edges depend on disjoint pairs of
variables, hence are independent.
(iii) \(X_0 \ge s\) and \(X_1 \ge t\) hold simultaneously iff
\(s + t \le P\), in which case each of the two variables has
\(P-s-t+1\) admissible values; thus
\(\E[X_0X_1] = \sum_{s,t\ge1} \Pr(X_0\ge s, X_1\ge t)
= \frac{1}{(P+1)^2} \sum_{r=1}^{P-1} (P-r) r^2
= \frac{P^2(P-1)}{12(P+1)}\), and subtracting \(\mu_P^2\) gives the
covariance.
\end{proof}

\begin{theorem}[variance of the hull dimension]\label{thm:C6}
\(\displaystyle
\operatorname{Var}\big(\dim_{\Fq} \Hull_k(C)\big)
= \sum_{s} \sum_{O \in \mathcal{O}_s} w_O^2\, V_{L_O}(P)\), where
\[
\begin{aligned}
V_1(P) &\text{ is as in Lemma~\ref{lem:varpieces}(i)},\\
V_2(P) &= \frac{P(P+2)(P^2+2P+3)}{18(P+1)^2},\\
V_L(P) &= \frac{L\, P(P+2)(P^2+2P+2)}{60(P+1)^2} \qquad (L \ge 3).
\end{aligned}
\]
\end{theorem}

\begin{proof}
Independence across orbits gives the weighted sum. For one orbit:
\(L=1\) is Lemma~\ref{lem:varpieces}(i). For \(L=2\),
\(\operatorname{Var}(S_2) = 2\operatorname{Var}(X) + 2\operatorname{Cov}(X_0,X_1)\);
with \(\operatorname{Var}(X) = \frac{P(P+2)(2P^2+4P+3)}{36(P+1)^2}\)
(from Lemma~\ref{lem:edge}) and (iii), the bracket factors as
\(P(P+2)(P^2+2P+3)\), giving \(V_2(P)\). For \(L \ge 3\),
\(\operatorname{Var}(S_L) = L\operatorname{Var}(X) + 2L\kappa_P\) (exactly
the \(L\) adjacent pairs contribute, by (ii)), and
\(\operatorname{Var}(X) + 2\kappa_P =
\frac{P(P+2)(P^2+2P+2)}{60(P+1)^2}\). At \(P=1\): \(V_1(1) = 0\),
\(V_2(1) = 1/4\), \(V_L(1) = L/16\). The \(L=2\) formula is not the
\(L \ge 3\) formula evaluated at \(L=2\) (at \(P=2\): \(44/81\) versus
\(8/27\)), because the two edges of a 2-orbit share both variables.
\end{proof}

\begin{corollary}[second-order structure of the two dimensions]\label{cor:secondorder}
Under the uniform model of this section,
\[
\operatorname{Var}\big(\dim_{\Fq} C\big)
= \frac{P(P+2)}{12}\sum_s\sum_{f\in\mathcal{F}_s} w_f^2,
\qquad
\operatorname{Cov}\big(\dim_{\Fq} C,\,\dim_{\Fq}\Hull_k(C)\big)=0.
\]
In particular, the two dimensions are uncorrelated (they need not be
independent).
\end{corollary}

\begin{proof}
By Theorem~\ref{thm:T4},
\(\dim_{\Fq} C=\sum_{s,f} w_f(P-A_f)\) with independent \(A_f\) uniform
on \(\{0,\dots,P\}\), and
\(\operatorname{Var}(P-A_f)=P(P+2)/12\), giving the variance formula.
For the covariance, orbit blocks are independent, so it suffices to
consider one orbit with \(K_O=\sum_i(P-A_i)\) and \(S_O=\sum_j X_j\).
In \(\operatorname{Cov}(K_O,S_O)\), a term
\(\operatorname{Cov}(P-A_i,X_j)\) can be nonzero only for
\(i\in\{j,j-1\}\). With
\(\gamma_P:=\operatorname{Cov}(B,\min(B,C))=P(P+2)/24\) for independent
\(B,C\) uniform on \(\{0,\dots,P\}\) (a direct power-sum computation),
the terms \(i=j\) and \(i=j-1\) contribute \(+\gamma_P\) and
\(-\gamma_P\) respectively when \(L\ge2\); for \(L=1\), symmetry under
\(A\mapsto P-A\) gives covariance \(0\) directly. Hence each orbit
contributes \(0\), and summing with the weights \(w_O^2\) gives
\(\operatorname{Cov}=0\).
\end{proof}

% =====================================================================
\section{Examples and computational verification}\label{sec:examples}
% =====================================================================

The six examples are arranged so that each one makes a different part
of the theory visible. Example~\ref{ex:E1} is small enough to check by
hand and shows how a single exponent tuple produces its two coordinates.
Example~\ref{ex:E2} isolates the simple-root specialization.
Examples~\ref{ex:E3} and \ref{ex:E4} leave the orbit model and compare
the exponent formulas with ranks of Gram matrices of explicitly
constructed codes, first for mixed orbit lengths and then for a long
orbit. Example~\ref{ex:E5} tests the product assembly over two
components, and Example~\ref{ex:E6} leaves the cyclic case \(\lambda = 1\)
used in Examples~\ref{ex:E3}--\ref{ex:E5} and works with a compatible
twist of order four. Throughout, \(K\) denotes the code-dimension coordinate and
\(H\) the hull-dimension coordinate of a labeled assignment.

\begin{example}[one orbit, \(P=2\), \(L=2\), \(w=1\), by hand]
\label{ex:E1}
This example shows the bookkeeping of Theorem~\ref{thm:T4} in the
smallest repeated-root case. The two positions form one orbit, and for
\(\alpha = (\alpha_0,\alpha_1)\) the two coordinates are
\(K = \sum_i (2-\alpha_i)\) and
\(H = \sum_i \min(2-\alpha_i, \alpha_{i-1})\). For instance, for
\(\alpha = (0,1)\) one gets \(K = (2-0) + (2-1) = 3\) and
\(H = \min(2,1) + \min(1,0) = 1\); the first edge contributes because
the predecessor exponent is positive, while the second edge contributes
nothing because its predecessor exponent is \(0\). Evaluating all nine
tuples gives
\[
\begin{aligned}
(0,0)&\mapsto(4,0),\quad (0,1)\mapsto(3,1),\quad (0,2)\mapsto(2,2),\quad
(1,0)\mapsto(3,1),\quad (1,1)\mapsto(2,2),\\
(1,2)&\mapsto(1,1),\quad (2,0)\mapsto(2,2),\quad (2,1)\mapsto(1,1),\quad (2,2)\mapsto(0,0).
\end{aligned}
\]
Grouping equal pairs yields the joint distribution
\(\{(4,0):1, (3,1):2, (2,2):3, (1,1):2, (0,0):1\}\). Two features are
already visible. First, the multiplicities are symmetric under
\(K \leftrightarrow 4-K\) at fixed \(H\), as
Corollary~\ref{cor:jointsym} requires for this orbit (here
\(\Delta = PwL = 4\)). Second, the moments can be read off directly:
\(\E[S_2] = (0\cdot1 + 1\cdot2 + 2\cdot3 + 1\cdot2 + 0\cdot1)/9 =
10/9 = 2\mu_2\), and
\(\operatorname{Var}(S_2) = 44/81 = V_2(2)\), the \(L=2\) case of
Theorem~\ref{thm:C6}. The separate \(L=2\) formula is essential here:
the two edges of a 2-orbit share both variables, and formally inserting
\(L=2\) into the \(L \ge 3\) formula would give the incorrect value
\(8/27\). Finally, exactly two assignments, \(\alpha \equiv 0\) and
\(\alpha \equiv P\), have hull contribution \(0\), as
Corollary~\ref{cor:C3} requires per orbit.
\end{example}

\begin{example}[simple-root orbit, \(P=1\), \(L=3\), \(w=1\)]\label{ex:E2}
At \(P=1\) the exponents are binary and the example checks the
simple-root specialization of Proposition~\ref{prop:P1}. The code
coordinate is the number of zeros, \(K = \sum_i (1-\alpha_i)\), while
the hull coordinate counts the cyclic transitions \(1 \to 0\). The
two constant words behave exceptionally: \(000\) gives \((K,H) =
(3,0)\) and \(111\) gives \((0,0)\). Every non-constant binary word
of length \(3\), viewed cyclically, has exactly one block of ones and
one block of zeros, hence exactly one \(1 \to 0\) transition and
\(H = 1\). There are three words of weight \(1\) (so \(K=2\)) and
three of weight \(2\) (so \(K=1\)). Therefore the joint distribution
is
\(\{(0,0):1,\ (1,1):3,\ (2,1):3,\ (3,0):1\}\), from which
\(\E = 3/4 = 3\mu_1\) and
\(\operatorname{Var} = 3/16 = V_3(1)\) follow immediately. The table
also displays the symmetry of Corollary~\ref{cor:jointsym} in its
smallest form: the hull multiplicity at \(K=0\) equals that at
\(K=3\), and the multiplicity at \(K=1\) equals that at \(K=2\).
\end{example}

\begin{example}[a repeated-root field case]\label{ex:E3}
Take \(K = \mathbb{F}_4\), \(k=1\), \(\lambda=1\), \(P=2\), \(n_0=5\),
\(n=10\). This case is chosen because it is genuinely repeated-root
(\(P=2\)) while still small, and because it contains both orbit types.
The factors of \(x^5-1\) over \(\mathbb{F}_4\) have degrees
\(1,2,2\); the degree-one factor is fixed by \(\tau\), and the two
degree-two factors form one orbit of length \(2\). Thus there are three
exponent positions and \(3^3=27\) assignments. (Dimensions in the
display are over \(K\) itself; the orbit weights are therefore \(1\)
for the length-one orbit and \(2\) for the length-two orbit.)

The point of the example is that the distribution below was computed
in two independent ways: from the exponent formulas of
Theorems~\ref{thm:T4} and \ref{thm:T6}, and directly from codes, by
forming, for each assignment \(a\), the generator \(g_a\), a basis of
shifts, and the rank of the code-first Gram matrix
\(\big(\sum_j B_{t,j}\sigma(B_{r,j})\big)\). The two computations
agree on the joint \((\dim_K C, \dim_K \Hull_1(C))\) distribution
\[
\begin{aligned}
\{&(0,0):1, (1,1):1, (2,0):1, (2,2):2, (3,3):2, (4,2):2, (4,4):3,
(5,5):3, (6,2):2,\\
&(6,4):3, (7,3):2, (8,0):1, (8,2):2, (9,1):1, (10,0):1\},
\end{aligned}
\]
whose entries sum to \(27\), as they must. The moment formulas can be
understood orbit by orbit: the length-one orbit is a self-loop and
contributes \(\nu_2 = 1/3\), whereas the length-two orbit has weight
\(2\) and two edges, hence contributes \(2 \cdot 2\mu_2 = 20/9\).
This gives
\(\E[\Hull] = 23/9 = \nu_2 + 2\cdot 2\mu_2 = \tfrac13 +
\tfrac{20}{9}\) --- the separate \(\nu_P\) term is not a cosmetic
distinction here, since using \(\mu_2\) for the length-one orbit as
well would give the wrong mean. The variance is
\(\operatorname{Var}(\Hull) = 194/81\) from Theorem~\ref{thm:C6}.
The counting corollaries are equally visible: there are two orbits, so
the LCD count is \(4 = 2^2\) (Corollary~\ref{cor:C3}), and the
self-orthogonal and dual-containing counts agree at \(12\), in line
with Corollary~\ref{cor:C4} and its refinement in
Corollary~\ref{cor:restricted}.
\end{example}

\begin{example}[orbit of length \(6\)]\label{ex:E4}
Take \(K = \mathbb{F}_8\), \(k=1\), \(\lambda=1\), \(P=2\), \(n_0=7\),
\(n=14\). The seven linear factors of \(x^7-1\) form \(\tau\)-orbits
of lengths \(1\) and \(6\), so there are \(3^7 = 2187\) assignments.
Whereas Example~\ref{ex:E3} tests mixed orbit lengths, this example is
included for the length-six orbit. The Gram-rank computation was
compared with the exponent formula assignment by assignment, not only
through the aggregated distribution; it gives
\(\E[\Hull] = 11/3\) and
\(\operatorname{Var}(\Hull) = 10/9 = V_1(2) + V_6(2) = \tfrac29 +
\tfrac89\), with LCD count \(4\) and self-orthogonal \(=\) dual-containing
count \(258\). The mean decomposition is again instructive: the
length-one orbit contributes \(\nu_2 = 1/3\), and the length-six orbit,
of weight \(1\) with six edges, contributes \(6\mu_2 = 10/3\).

For the record, the full joint distribution of
\((\dim_{\Fq} C, \dim_{\Fq}\Hull_k(C))\) is
{\footnotesize
\((0,0){:}\,1\), \((1,1){:}\,7\), \((2,0){:}\,1\), \((2,2){:}\,27\), \((3,1){:}\,6,
(3,2){:}\,12\), \((3,3){:}\,59\), \((4,2){:}\,33\), \((4,3){:}\,48\), \((4,4){:}\,80,
(5,2){:}\,24\), \((5,3){:}\,86\), \((5,4){:}\,96\), \((5,5){:}\,60\), \((6,2){:}\,24,
(6,3){:}\,84\), \((6,4){:}\,144\), \((6,5){:}\,84\), \((6,6){:}\,21\), \((7,2){:}\,24,
(7,3){:}\,84\), \((7,4){:}\,156\), \((7,5){:}\,102\), \((7,6){:}\,24\), \((7,7){:}\,3,
(8,2){:}\,24\), \((8,3){:}\,84\), \((8,4){:}\,144\), \((8,5){:}\,84\), \((8,6){:}\,21,
(9,2){:}\,24\), \((9,3){:}\,86\), \((9,4){:}\,96\), \((9,5){:}\,60\), \((10,2){:}\,33,
(10,3){:}\,48\), \((10,4){:}\,80\), \((11,1){:}\,6\), \((11,2){:}\,12,
(11,3){:}\,59\), \((12,0){:}\,1\), \((12,2){:}\,27\), \((13,1){:}\,7\), \((14,0){:}\,1\).}

The exponent formula is thus checked against the hull dimensions of
actual codes on an orbit of length \(6\). One point must be stated
correctly: a dimension check of this kind cannot certify the predecessor
orientation, because the hull dimension does not depend on the
orientation convention at all --- for all \(a, b\),
\(\min(P-a,b) - \min(P-b,a) = b - a\), which telescopes to \(0\) around
any orbit, so the predecessor and successor readings give the same hull
dimension for every assignment (cf.\ the remark after
Theorem~\ref{thm:T6}). The orientation is instead visible in the hull
generator: here, the exponent of \(f_i\) in \(h^{\#}\) equals the
predecessor value \(P - \alpha_{i-1}\) for every one of the \(2187\)
assignments, while the successor value \(P - \alpha_{i+1}\) agrees only
for \(27\) of them (checked in
\texttt{code/validate\_audit\_fixes.py}). The predecessor orientation is
fixed by the code-first pairing (Theorem~\ref{thm:T3}).
\end{example}

\begin{example}[two components]\label{ex:E5}
The final example tests the assembly step
(Proposition~\ref{prop:assembly} and Theorem~\ref{thm:T6}), rather
than a single orbit. Let \(A = \mathbb{F}_2 \times \mathbb{F}_4\)
(\(q=2\), \(k=0\), \(\lambda=(1,1)\), \(P=2\), \(n_0=7\)). The orbit
data \((L_O, w_O)\) are \((1,1), (2,3)\) on the first component and
\((1,2), (2,6)\) on the second: four orbits and
\(1+2+1+2 = 6\) exponent positions in total, hence
\(3^6 = 729\) labeled assignments. Because assignments on distinct
orbits are independent, the global joint distribution is the
convolution (equivalently, the coefficientwise product) of the four
per-orbit distributions of Theorem~\ref{thm:T5}; the same distribution
is obtained by computing the two components separately and adding their
dimensions. This agreement is the concrete content of the product
formula in Theorem~\ref{thm:T6}.

The summary numbers illustrate how the global formulas are assembled.
There are four orbits, so the LCD count is \(2^4 = 16\)
(Corollary~\ref{cor:C3}). The self-orthogonal and dual-containing
counts are \(144\) each, as Corollary~\ref{cor:C4} predicts. The mean
splits according to orbit length: the two length-one orbits contribute
\((1+2)\nu_2 = 1\), while the two length-two orbits contribute
\((3\cdot 2 + 6\cdot 2)\mu_2 = 10\), giving
\(\E[\Hull] = 11\) (Theorem~\ref{thm:C5}). Likewise the variance is
the weight-squared sum
\(\operatorname{Var}(\Hull) = 230/9\) of
Theorem~\ref{thm:C6}, namely
\(1^2 V_1(2) + 3^2 V_2(2) + 2^2 V_1(2) + 6^2 V_2(2) = 230/9\).
Thus Example~\ref{ex:E5} verifies, in one place, the independence of
orbit blocks, the \(L=1\) versus \(L \ge 2\) mean constants, and the
the \(w_O^2\) weighting in the variance.
\end{example}

\begin{example}[a twist of order four]\label{ex:E6}
All field-level examples so far have \(\lambda = 1\). Take instead
\(K = \mathbb{F}_9 = \mathbb{F}_3(\iota)\) with \(\iota^2 = -1\),
\(k = 1\), \(\lambda = \iota\) of order \(4\) (compatible, since
\(\lambda^{1+3} = \lambda^4 = 1\)), \(P = 3\), \(n_0 = 4\), \(n = 12\).
Here \(\mu = -\iota\), and \(x^4 - \mu = x^4 + \iota\) factors over
\(\mathbb{F}_9\) as two irreducible quadratics forming a single
\(\tau\)-orbit of length \(2\); hence there are \((3+1)^2 = 16\)
assignments. The joint distribution of
\((\dim_{\Fq} C, \dim_{\Fq}\Hull_k(C))\) is
\(\{(0,0):1, (2,2):2, (4,4):3, (6,6):4, (8,4):3, (10,2):2, (12,0):1\}\);
hence \(\E[\Hull] = 7/2\) and \(\operatorname{Var}(\Hull) = 15/4\), the
LCD count is \(2 = 2^1\) (one orbit), and the self-orthogonal and
dual-containing counts are \(10\) each. Every value was recomputed by
Gram-rank brute force over the explicitly constructed codes
(\texttt{code/validate\_audit\_fixes.py}), so the constacyclic theory is
exercised here beyond the cyclic case \(\lambda = 1\) and the negacyclic
case \(\lambda = -1\).
\end{example}

\paragraph{A table of EAQECC-optimal and near-optimal codes.}
We record ten further finite-field codes, all with \(\lambda = 1\) (hence
\(\mu = 1\)) and a single component (\(N = m_1 = 1\)), so all dimensions
below are over \(\mathbb{F}_q\) itself. Throughout, \(\omega\) denotes a
root of \(x^2 + x + 1\) in \(\mathbb{F}_4\), \(\theta\) a root of
\(x^3 + x + 1\) in \(\mathbb{F}_8\), and \(\iota\) an element of
\(\mathbb{F}_9\) with \(\iota^2 = -1\). The exponent tuple \(a\) refers to
the following ordered factor lists of \(x^{n_0} - 1\): over
\(\mathbb{F}_2\) with \(n_0 = 3\): \((x+1,\ x^2+x+1)\); over
\(\mathbb{F}_3\) with \(n_0 = 2\): \((x+1,\ x+2)\); over \(\mathbb{F}_4\)
with \(n_0 = 5\): \((x+1,\ x^2 + \omega x + 1,\ x^2 + \omega^2 x + 1)\);
over \(\mathbb{F}_4\) with \(n_0 = 3\):
\((x+1,\ x+\omega,\ x+\omega^2)\); over \(\mathbb{F}_5\) with \(n_0 = 2\):
\((x+1,\ x+4)\); over \(\mathbb{F}_7\) with \(n_0 = 2\): \((x+1,\ x+6)\);
over \(\mathbb{F}_8\) with \(n_0 = 7\): \((x+\beta)_\beta\) with \(\beta\)
running through \(1, \theta, 1+\theta, \theta^2, 1+\theta^2,
\theta+\theta^2, 1+\theta+\theta^2\); over \(\mathbb{F}_8\) with
\(n_0 = 3\): \((x+1,\ x^2+x+1)\); over \(\mathbb{F}_9\) with \(n_0 = 8\):
\((x+\beta)_\beta\) with \(\beta\) running through \(1, 2, \iota,
1+\iota, 2+\iota, 2\iota, 1+2\iota, 2+2\iota\); and over
\(\mathbb{F}_9\) with \(n_0 = 2\): \((x+1,\ x+2)\). For each row, \(d\)
is the minimum distance of \(C(a)\), determined independently of the
enumerator (exact computation; see the computational-verification
paragraph below), and \(H = \dim \Hull_k(C(a))\) is computed from
Theorem~\ref{thm:T4}. The EAQECC column applies the conditional quantum
remark of Section~\ref{sec:limitations}, giving
\([[n, K-H, d;\, n-K-H]]_q\). The last column is the EA Singleton defect
\(\Delta = \big(n - (K-H) + (n-K-H)\big) - 2(d-1) = 2(n-K-d+1)\).
Following the convention of~\cite{DIMP2026}, we write \({}^*\) when
\(\Delta = 0\), i.e.\ the EAQECC meets the EA Singleton bound (we call
such a code EA-optimal), and \({}^\dagger\) when \(\Delta = 2\)
(near-optimal). Optimality here is meant only with respect to this EA
Singleton bound; no comparison with best-known classical code tables is
made. Rows with \(P = 1\) are simple-root specializations; the remaining
rows are genuinely repeated-root.

\begin{table}[htbp]
\centering
\footnotesize
\caption{Ten new finite-field exponent assignments. The EAQECC column
gives the optimal code \([[n, K-H, d;\, n-K-H]]_q\) obtained from the
conditional quantum remark, marked \({}^*\) if it meets the EA Singleton
bound (\(\Delta = 0\)) and \({}^\dagger\) if \(\Delta = 2\).}
\label{tab:eaqecc}
\begin{tabular}{cllcccr}
\hline
No. & Setting & \(a\) & \([n,K,d]_q\) & \(H\) & EAQECC & \(\Delta\) \\
\hline
1 & \(\mathbb{F}_2\ (k{=}0, P{=}2, n_0{=}3)\) & \((2,0)\) & \([6,4,2]_2\) & \(0\) & \([[6,4,2;2]]_2^{\dagger}\) & \(2\) \\
2 & \(\mathbb{F}_3\ (k{=}0, P{=}3, n_0{=}2)\) & \((1,3)\) & \([6,2,4]_3\) & \(1\) & \([[6,1,4;3]]_3^{\dagger}\) & \(2\) \\
3 & \(\mathbb{F}_4\ (k{=}1, P{=}1, n_0{=}5)\) & \((0,0,1)\) & \([5,3,3]_4\) & \(2\) & \([[5,1,3;0]]_4^{*}\) & \(0\) \\
4 & \(\mathbb{F}_4\ (k{=}1, P{=}2, n_0{=}3)\) & \((0,1,2)\) & \([6,3,3]_4\) & \(1\) & \([[6,2,3;2]]_4^{\dagger}\) & \(2\) \\
5 & \(\mathbb{F}_5\ (k{=}0, P{=}5, n_0{=}2)\) & \((1,3)\) & \([10,6,4]_5\) & \(3\) & \([[10,3,4;1]]_5^{\dagger}\) & \(2\) \\
6 & \(\mathbb{F}_7\ (k{=}0, P{=}7, n_0{=}2)\) & \((1,3)\) & \([14,10,4]_7\) & \(4\) & \([[14,6,4;0]]_7^{\dagger}\) & \(2\) \\
7 & \(\mathbb{F}_8\ (k{=}1, P{=}1, n_0{=}7)\) & \((0,0,0,0,1,1,1)\) & \([7,4,4]_8\) & \(1\) & \([[7,3,4;2]]_8^{*}\) & \(0\) \\
8 & \(\mathbb{F}_8\ (k{=}1, P{=}2, n_0{=}3)\) & \((2,0)\) & \([6,4,2]_8\) & \(0\) & \([[6,4,2;2]]_8^{\dagger}\) & \(2\) \\
9 & \(\mathbb{F}_9\ (k{=}1, P{=}1, n_0{=}8)\) & \((0,1,0,0,0,1,1,1)\) & \([8,4,5]_9\) & \(2\) & \([[8,2,5;2]]_9^{*}\) & \(0\) \\
10 & \(\mathbb{F}_9\ (k{=}1, P{=}3, n_0{=}2)\) & \((1,2)\) & \([6,3,3]_9\) & \(2\) & \([[6,1,3;1]]_9^{\dagger}\) & \(2\) \\
\hline
\end{tabular}
\end{table}

\paragraph{Computational verification.}
All formulas above were checked by exact (integer and rational)
computation; the scripts and their saved outputs accompany the paper in
the directory \texttt{code/}. (i) \emph{Orbit level}
(\texttt{code/validate\_p1\_14.py}): direct enumeration of exponent
tuples against the transfer traces of Theorem~\ref{thm:T5}, on the grid
of \(52\) pairs \((P,L)\) with \(P \in \{1,2,3,4,5,7,8,9\}\),
\(1 \le L \le 8\) and \((P+1)^L \le 10^5\); the moment closed forms were
also re-derived symbolically from their defining sums. (ii) \emph{Field
level}: Gram-rank brute force over explicitly constructed codes for
\(\mathbb{F}_4\) with \(n_0 = 3\) (\(27\) assignments) and \(n_0 = 5\)
(\(27\)), \(\mathbb{F}_9\) with \(k = 1\), \(\lambda = 1\) (\(256\)) and
\(\lambda = -1\) (\(256\)), \(\mathbb{F}_3\) with \(\lambda = -1\)
(\(16\)), and \(\mathbb{F}_8\) with \(n_0 = 7\) (\(2187\)) ---
\(2769\) field-level assignments in total --- together with the
two-component assembly of Example~\ref{ex:E5} (\(729\) assignments),
checked by convolution of the component distributions. (iii) The twist
of order four in Example~\ref{ex:E6} (\(16\) assignments) and further
field-level cases with \(P \in \{4, 8, 9\}\) (\(\mathbb{F}_2\) with
\(P = 4\) (\(25\)) and \(P = 8\) (\(81\)); \(\mathbb{F}_3\) with
\(P = 9\) (\(100\))) are checked in
\texttt{code/validate\_audit\_fixes.py}, which also verifies the
hull-generator orientation statement of Example~\ref{ex:E4}. (iv) The
ten codes of Table~\ref{tab:eaqecc} were checked separately by the
accompanying script \texttt{code/validate\_code\_table.py}: for every
row the hull dimension was computed both from Theorem~\ref{thm:T4} and
from the code-first Gram matrix, and the minimum distance exactly by
parity-check column dependency, confirmed by an explicit minimum-weight
codeword and, except in the \(\mathbb{F}_7\) row, by direct enumeration
of all codewords. These computations are supporting evidence for the
tested cases; the proofs of
Sections~\ref{sec:param}--\ref{sec:moments} are independent of them.

% =====================================================================
\section{Limitations and scope}\label{sec:limitations}
% =====================================================================

\emph{Scope.} All enumerations count labeled exponent assignments over
a fixed labeled decomposition \(A \cong \prod_s K_s\); they are not
counts of equivalence or isometry classes of codes, and no such
interpretation should be attached to \(\mathcal{E}(u,z)\). Only
compatible twists are treated: for incompatible \(\lambda_s\),
Theorem~\ref{thm:T3} persists in the two-modulus form
\(x^n - \lambda_s^{-p^{d_s-k}}\), but the hull formula and everything
following it are not asserted. The alphabet is reduced (square-free) by
the standing hypothesis \(A \cong \prod_s K_s\); non-reduced alphabets,
such as chain rings, are outside the scope, as are distances, which the
enumerator does not address. The means and variances of
Section~\ref{sec:moments} refer to the uniform distribution on labeled
assignments; orbit contributions are independent in this model, edges
within an orbit are not (Lemma~\ref{lem:varpieces}), and the \(L=1\),
\(L=2\), \(L\ge3\) regimes are genuinely distinct. Finally, the
predecessor index \(\alpha_{i-1}\) in
Theorems~\ref{thm:T3}--\ref{thm:T6} is tied to the code-first pairing of
Section~\ref{sec:prelim}: reversing the convention replaces \(\rho\) by
\(\sigma\) and \(\alpha_{i-1}\) by \(\alpha_{i+1}\). By the telescoping
identity \(\min(P-a,b) - \min(P-b,a) = b - a\), every hull dimension ---
and hence every statistic of Sections~\ref{sec:corollaries}
and~\ref{sec:moments} --- is unchanged assignment by assignment; the
convention is visible only in the identification of the hull generator
(Theorem~\ref{thm:T4}), not in any dimension.

\emph{Open problems.} Natural directions include the incompatible-twist
regime, where the dual exists in a two-modulus form but the hull
intersection must be reconsidered; counting up to monomial equivalence
rather than by labeled assignments; the extension from constacyclic
(index-one) to multi-twisted codes; and distance distributions, which
would require machinery beyond the enumerator.

\begin{remark}[conditional quantum remark]
In the finite-field case (\(N=1\), \(m_1=1\), so \(q = p^e\) and
\(0 \le k < e\)), let \(C\) be an \([n, K]\) code over \(\Fq\) and let
\(H = \dim \Hull_k(C)\). Let
\(\Pi\) be a parity-check matrix of \(C\) (rows a basis of \(C^\perp\)) and
put \(c = \mathrm{rank}\big(\Pi\, \rho(\Pi)^{\mathsf T}\big)\), with
\(\rho\) applied entrywise. The kernel of the pairing map
\(D(C) \to (C^\perp)^*\), \(y \mapsto \big(z \mapsto \sum_i z_i y_i\big)\),
is exactly \(C \cap D(C)\); hence \(c = (n-K) - H\). Therefore, when the
entanglement-assisted construction is used in its Galois form
\cite{DIMP2026} (the framework of \cite{BDS2006}), \(C\) yields an
\([[n, K-H, d_{\mathrm{EA}};\, n-K-H]]\) entanglement-assisted quantum
code, where \(d_{\mathrm{EA}}\) is the distance furnished by that
construction and must be determined independently. The enumerator of
Theorem~\ref{thm:T6} determines the joint distribution of \((K, H)\)
only; it yields no distance information, and no quantum parameter is
claimed from it alone. The product case \(N \ge 2\) requires, in
addition, control of distances under the product decomposition. The ten
codes of Table~\ref{tab:eaqecc} are explicit instances whose distances
were computed independently, so their displayed parameters are
unconditional; no asymptotic family is claimed, since for repeated-root
families the distance --- which the enumerator does not control --- is
the limiting quantity.
\end{remark}

% =====================================================================
\section{Conclusion}
% =====================================================================

For repeated-root constacyclic codes over a finite product of finite
fields with compatible twists, the joint distribution of code dimension and
Galois hull dimension over all labeled exponent assignments is given by
a product of transfer-matrix traces indexed by the orbits of the
reciprocal involution. The enumerator yields the total count, the
code-dimension distribution, the LCD count \(2^{\#\mathrm{orbits}}\), the
self-orthogonal and dual-containing counts, and exact mean and variance
formulas in which length-one orbits enter through the separate constant
\(\nu_P\). At \(P=1\) the theory reduces to the binary simple-root
enumerator. Natural directions for further work include the
incompatible-twist regime, where the dual already exists in a
two-modulus form but the hull intersection must be reconsidered;
counting up to monomial equivalence rather than by labeled assignments;
and distance questions, which the enumerator alone does not address.

% =====================================================================
\begin{thebibliography}{99}
% =====================================================================

\bibitem{ACLMRS2024}
S.~E. Anderson, E.~Camps-Moreno, H.~H. L\'opez, G.~L. Matthews,
D.~Ruano, and I.~Soprunov,
Relative hulls and quantum codes,
\emph{IEEE Transactions on Information Theory} 70(5) (2024),
3190--3201. doi:10.1109/TIT.2024.3373550.

\bibitem{BBO2026}
S.~Bouyuklieva, I.~Bouyukliev, and F.~\"Ozbudak,
Sequence of numbers of linear codes with increasing hull dimensions,
\emph{Designs, Codes and Cryptography} 94 (2026), Article 39.
doi:10.1007/s10623-025-01776-9.

\bibitem{BDS2006}
T.~Brun, I.~Devetak, and M.-H.~Hsieh,
Correcting quantum errors with entanglement,
\emph{Science} 314 (2006), 436--439. doi:10.1126/science.1131563.

\bibitem{CDL2015}
B.~Chen, H.~Q. Dinh, and H.~Liu,
Repeated-root constacyclic codes of length \(2\ell^m p^n\),
\emph{Finite Fields and Their Applications} 33 (2015), 137--159.
doi:10.1016/j.ffa.2014.11.006.

\bibitem{DP2025}
I.~Debnath and O.~Prakash,
Average dimensions of Galois hulls of constacyclic codes,
\emph{Advances in Mathematics of Communications} 19(6) (2025),
1569--1604. doi:10.3934/amc.2025010.

\bibitem{DPDC2025}
I.~Debnath and O.~Prakash,
Hulls of separable double cyclic codes over \(\mathbb{Z}_{p^2}\),
\emph{Computational and Applied Mathematics} 44 (2025), Article 255.
doi:10.1007/s40314-025-03206-x.

\bibitem{DIMP2026}
I.~Debnath, H.~Islam, E.~Mart\'inez-Moro, and O.~Prakash,
Galois hulls of constacyclic codes over affine algebra rings,
\emph{Discrete Mathematics} 349(2) (2026), Article 114750.
doi:10.1016/j.disc.2025.114750; arXiv:2412.08512.

\bibitem{DIYP2026}
I.~Debnath, H.~Islam, S.~Yadav, and O.~Prakash,
Study of small Galois hull dimensions of constacyclic codes,
\emph{Advances in Mathematics of Communications} 22 (2026), 1--18.
doi:10.3934/amc.2025054.

\bibitem{DPI2023}
I.~Debnath, O.~Prakash, and H.~Islam,
Galois hulls of constacyclic codes over finite fields,
\emph{Cryptography and Communications} 15(1) (2023), 111--127.
doi:10.1007/s12095-022-00591-6.

\bibitem{DPI2023corr}
I.~Debnath, O.~Prakash, and H.~Islam,
Correction to: Galois hulls of constacyclic codes over finite fields,
\emph{Cryptography and Communications} 15 (2023), 129--130.
doi:10.1007/s12095-022-00602-6.

\bibitem{DS2024}
S.~T. Dougherty and S.~\c{S}ahinkaya,
On cyclic and negacyclic codes with one-dimensional hulls and their
applications,
\emph{Advances in Mathematics of Communications} 18(5) (2024),
1364--1378. doi:10.3934/amc.2022096.

\bibitem{ES2026}
R.~F.~Taki Eldin and P.~Sol\'e,
Generator polynomial matrices of the Galois hulls of multi-twisted codes,
\emph{Finite Fields and Their Applications} 109 (2026), Article 102712.
doi:10.1016/j.ffa.2025.102712; arXiv:2305.15237.

\bibitem{FZ2017}
Y.~Fan and L.~Zhang,
Galois self-dual constacyclic codes,
\emph{Designs, Codes and Cryptography} 84 (2017), 473--492.
doi:10.1007/s10623-016-0282-8.

\bibitem{HP2003}
W.~C. Huffman and V.~Pless,
\emph{Fundamentals of Error-Correcting Codes},
Cambridge University Press, 2003.

\bibitem{LP2020}
H.~Liu and X.~Pan,
Galois hulls of linear codes over finite fields,
\emph{Designs, Codes and Cryptography} 88 (2020), 241--255.
doi:10.1007/s10623-019-00681-2.

\bibitem{LS2024}
S.~Li and M.~Shi,
Characterization and classification of binary linear codes with various
hull dimensions from an improved mass formula,
\emph{IEEE Transactions on Information Theory} 70(5) (2024),
3357--3372. doi:10.1109/TIT.2023.3298379.

\bibitem{LSL2025}
S.~Li, M.~Shi, and S.~Ling,
A mass formula for linear codes with prescribed hull dimension and
related classification,
\emph{IEEE Transactions on Information Theory} 71(1) (2025),
273--286. doi:10.1109/TIT.2024.3425058.

\bibitem{LYW2025}
Y.~Liu, X.~Yu, and G.~Wang,
Several constructions of linear codes with various Galois hull
dimensions,
\emph{Computational and Applied Mathematics} 44 (2025), Article 299.
doi:10.1007/s40314-025-03261-4.

\bibitem{MS1977}
F.~J. MacWilliams and N.~J.~A. Sloane,
\emph{The Theory of Error-Correcting Codes},
North-Holland, 1977.

\bibitem{Sendrier1997}
N.~Sendrier,
On the dimension of the hull,
\emph{SIAM Journal on Discrete Mathematics} 10(2) (1997), 282--293.
doi:10.1137/S0895480195294027.

\bibitem{Sharma2015}
A.~Sharma,
Repeated-root constacyclic codes of length \(\ell^t p^s\) and their
dual codes,
\emph{Cryptography and Communications} 7 (2015), 229--255.
doi:10.1007/s12095-014-0106-5.

\bibitem{Talbi2021}
S.~Talbi, A.~Batoul, A.~Fotue Tabue, and E.~Mart\'inez-Moro,
Galois hulls of cyclic serial codes over a finite chain ring,
arXiv:2102.06995 (2021).

\bibitem{WZLL2023}
T.~Wu, S.~Zhu, L.~Liu, and L.~Li,
Repeated-root constacyclic codes of length \(6\ell^m p^n\),
\emph{Advances in Mathematics of Communications} 17(5) (2023),
1154--1180. doi:10.3934/amc.2021044.

\bibitem{ZKZ2026}
E.~Zhang, B.~Kong, and X.~Zheng,
Quantum codes from Galois hulls of constacyclic codes over a finite
non-chain ring,
\emph{Entropy} 28 (2026), 407. doi:10.3390/e28040407.

\bibitem{ZWLL2025}
Q.~Zhang, W.~Wang, S.~Luo, and Y.~Li,
Repeated-root constacyclic codes of length \(k^s\ell^m p^n\) over
finite fields,
\emph{Finite Fields and Their Applications} 101 (2025), Article 102542.
doi:10.1016/j.ffa.2024.102542.

\end{thebibliography}

\end{document}
